% =========================================================================
% EXPERIMENT 5: SQL TABLE OPERATIONS, REFERENTIAL INTEGRITY, AND NESTED QUERIES
% =========================================================================
\label{exp:nested_queries_practice}
\label{exp:practical_05}
\noindent
\begin{center}
    {\LARGE \textbf{DBMS LABORATORY PRACTICAL}}\\[0.25cm]
    \textbf{Practical No.:} 5 \hfill \textbf{Date:} 26/10/2026 \hfill \textbf{Page No.:} \thepage
\end{center}
\vspace{-0.2cm}
\hrule
\vspace{0.35cm}

\subsection*{Question / Problem Statement}
Consider the database \texttt{CollegeLab} comprising relational tables \texttt{Students}, \texttt{ExamMarks}, \texttt{Department}, and \texttt{Employee}. Write, optimize, and execute SQL queries to perform table operations, enforce referential integrity, and evaluate subqueries:
\begin{enumerate}[noitemsep,topsep=2pt]
    \item[(i)] Display students from Guwahati in descending order of marks.
    \item[(ii)] Find the number of students in each course who scored at least 70.
    \item[(iii)] Display the second-highest distinct marks using a subquery.
    \item[(iv)] Find employees earning more than the average salary of all employees.
    \item[(v)] List students whose \texttt{Roll\_No} appears in \texttt{ExamMarks} with \texttt{Score} $\ge 90$.
    \item[(vi)] Find departments having no employees using \texttt{NOT EXISTS}.
    \item[(vii)] Display the highest-paid employee in each department using a correlated subquery.
    \item[(viii)] List students with scores in both DBMS and Python using nested queries.
    \item[(ix)] Create a \texttt{CourseDetails} table and establish a foreign key with \texttt{Students}.
    \item[(x)] Demonstrate and explain why deleting a parent row fails when dependent child records exist.
\end{enumerate}

\subsection*{1. Database Schema Setup \& Benchmark Data Insertion}

\subsubsection*{SQL Query}
\begin{lstlisting}[basicstyle=\ttfamily\footnotesize,caption={Creating CollegeLab schema and populating benchmark records}]
CREATE DATABASE IF NOT EXISTS CollegeLab;
USE CollegeLab;

-- 1. Create relational tables with primary and foreign key constraints
CREATE TABLE Students (
    Roll_No INT PRIMARY KEY,
    Name VARCHAR(50) NOT NULL,
    Address VARCHAR(60),
    Age INT,
    Course VARCHAR(50),
    Marks DECIMAL(5, 2)
);

CREATE TABLE ExamMarks (
    Mark_ID INT PRIMARY KEY,
    Student_ID INT NOT NULL,
    Subject VARCHAR(50) NOT NULL,
    Score DECIMAL(5, 2) NOT NULL,
    CONSTRAINT fk_exam_student FOREIGN KEY (Student_ID)
        REFERENCES Students(Roll_No)
);

CREATE TABLE Department (
    Dept_ID INT PRIMARY KEY,
    Dept_Name VARCHAR(50) NOT NULL
);

CREATE TABLE Employee (
    Emp_ID INT PRIMARY KEY,
    Emp_Name VARCHAR(50) NOT NULL,
    Salary DECIMAL(10, 2) NOT NULL,
    Dept_No INT,
    CONSTRAINT fk_employee_dept FOREIGN KEY (Dept_No)
        REFERENCES Department(Dept_ID)
);

-- 2. Populate sample records
INSERT INTO Students VALUES
(1, 'Rahul', 'Guwahati', 20, 'MCA', 85), (2, 'Priya', 'Jorhat', 21, 'MCA', 92),
(3, 'Amit', 'Dibrugarh', 22, 'BCA', 68), (4, 'Neha', 'Guwahati', 20, 'MCA', 78),
(5, 'Riya', 'Silchar', 23, 'BCA', 92), (6, 'Arjun', 'Tezpur', 21, 'MCA', 60);

INSERT INTO ExamMarks VALUES
(101, 1, 'DBMS', 88), (102, 1, 'Python', 82), (103, 2, 'DBMS', 95), (104, 2, 'Python', 89),
(105, 3, 'DBMS', 70), (106, 4, 'DBMS', 80), (107, 5, 'Python', 94), (108, 6, 'DBMS', 58);

INSERT INTO Department VALUES
(10, 'IT'), (20, 'HR'), (30, 'Finance'), (40, 'Research');

INSERT INTO Employee VALUES
(201, 'Anita', 55000, 10), (202, 'Bikash', 48000, 10),
(203, 'Farah', 62000, 20), (204, 'Deep', 45000, 30), (205, 'Kabir', 70000, 20);
\end{lstlisting}

\newpage

\subsection*{2. Filtering, Aggregations, and Nested Subqueries (Exercises i--v)}

\subsubsection*{SQL Query}
\begin{lstlisting}[basicstyle=\ttfamily\footnotesize,caption={Executing conditional filtering, grouping, scalar subqueries, and set membership}]
-- (i) Display students from Guwahati in descending order of marks
SELECT * FROM Students
WHERE Address = 'Guwahati'
ORDER BY Marks DESC;

-- (ii) Find number of students in each course scoring at least 70
SELECT Course, COUNT(*) AS StudentCount
FROM Students
WHERE Marks >= 70
GROUP BY Course;

-- (iii) Display second-highest distinct marks using a subquery
SELECT MAX(Marks) AS SecondHighestMarks
FROM Students
WHERE Marks < (SELECT MAX(Marks) FROM Students);

-- (iv) Employees earning more than average salary of all employees
SELECT Emp_ID, Emp_Name, Salary, Dept_No
FROM Employee
WHERE Salary > (SELECT AVG(Salary) FROM Employee);

-- (v) Students whose Roll_No appears in ExamMarks with Score >= 90
SELECT Roll_No, Name, Course, Marks
FROM Students
WHERE Roll_No IN (
    SELECT Student_ID FROM ExamMarks WHERE Score >= 90
);
\end{lstlisting}

\newpage

\subsection*{3. Correlated Subqueries, Constraints, and Referential Integrity (Exercises vi--x)}

\subsubsection*{SQL Query}
\begin{lstlisting}[basicstyle=\ttfamily\footnotesize,caption={Correlated subqueries, NOT EXISTS, multi-subqueries, and foreign keys}]
-- (vi) Find departments having no employees using NOT EXISTS
SELECT Dept_ID, Dept_Name
FROM Department d
WHERE NOT EXISTS (
    SELECT 1 FROM Employee e WHERE e.Dept_No = d.Dept_ID
);

-- (vii) Highest-paid employee in each department (correlated subquery)
SELECT e.Emp_ID, e.Emp_Name, e.Salary, e.Dept_No
FROM Employee e
WHERE e.Salary = (
    SELECT MAX(e2.Salary)
    FROM Employee e2
    WHERE e2.Dept_No = e.Dept_No
);

-- (viii) Students with scores in both DBMS and Python (nested subqueries)
SELECT Roll_No, Name, Course
FROM Students
WHERE Roll_No IN (
    SELECT Student_ID FROM ExamMarks WHERE Subject = 'DBMS'
)
AND Roll_No IN (
    SELECT Student_ID FROM ExamMarks WHERE Subject = 'Python'
);

-- (ix) Create CourseDetails table and establish foreign key relationship
CREATE TABLE CourseDetails (
    Course_Code VARCHAR(50) PRIMARY KEY,
    Duration_Years INT NOT NULL CHECK (Duration_Years > 0),
    Department VARCHAR(50) NOT NULL
);
INSERT INTO CourseDetails VALUES
('MCA', 2, 'Computer Science & Applications'),
('BCA', 3, 'Computer Applications');

ALTER TABLE Students
ADD CONSTRAINT fk_students_coursedetails
FOREIGN KEY (Course) REFERENCES CourseDetails(Course_Code);

-- (x) Demonstrate referential integrity violation on parent record deletion
DELETE FROM Students WHERE Roll_No = 1;
\end{lstlisting}

\vspace{0.25cm}
\subsection*{4. Referential Integrity Analysis (Exercise x)}
\noindent\textbf{Theoretical Explanation:}\quad In relational databases, referential integrity mandates that every foreign key value in a child table (\texttt{ExamMarks.Student\_ID}) must reference a valid, existing primary key in the parent table (\texttt{Students.Roll\_No}). When executing:
\begin{center}
    \texttt{DELETE FROM Students WHERE Roll\_No = 1;}
\end{center}
the RDBMS halts execution and returns \textbf{Error 1451 (23000)}: \textit{"Cannot delete or update a parent row: a foreign key constraint fails"}. Because child records in \texttt{ExamMarks} reference \texttt{Student\_ID = 1}, deleting the parent record would generate orphan child records. Under MySQL's default referential action (\texttt{RESTRICT}/\texttt{NO ACTION}), the delete operation is rejected. To delete successfully, either the referencing child rows must be removed first, or the foreign key must be defined with \texttt{ON DELETE CASCADE} or \texttt{ON DELETE SET NULL}.

\vspace{0.35cm}
\noindent\textbf{Result / Observation:}\quad Relational table operations, aggregate summaries, single-row subqueries, correlated subqueries, set membership checks (\texttt{IN}, \texttt{NOT EXISTS}), and referential integrity constraints were implemented and evaluated in MySQL according to the laboratory exercises. Constraint enforcement against orphan record creation and foreign key violation behavior were verified.
